\documentclass[letterpaper, 10 pt, conference]{ieeeconf}
\IEEEoverridecommandlockouts
\overrideIEEEmargins

\usepackage{amsmath}
\usepackage{amssymb}
\usepackage{booktabs}
\usepackage{graphicx}
\usepackage{url}

\title{\LARGE \bf
Model Compression for Deep Neural Networks:\\
Pruning and Quantization on the Fruits-360 Dataset
}

\author{
  \parbox{3in}{\centering
    Person 1$^{1}$, Person 2$^{1}$, Person 3$^{1}$, Person 4$^{1}$
    \thanks{$^{1}$Faculty of Engineering, University of Rijeka, Croatia}
  }
}

\begin{document}

\maketitle
\thispagestyle{empty}
\pagestyle{empty}

%%%%%%%%%%%%%%%%%%%%%%%%%%%%%%%%%%%%%%%%%%%%%%%%%%%%%%%%%%%%%%%%%%%%%%%%%%%%%%%%
\begin{abstract}
Deep neural networks (DNNs) have achieved remarkable results in image
classification tasks, yet their growing complexity makes deployment on
resource-constrained devices challenging.
This paper presents a practical implementation of two model compression
techniques --- weight pruning and post-training quantization --- applied to
a convolutional neural network (CNN) trained from scratch on the Fruits-360
dataset.
The entire pipeline is implemented in pure NumPy without any deep-learning
framework.
Our results show that removing 50\% of the lowest-magnitude weights
(global unstructured pruning) reduces the network's sparsity to 50\%
while incurring an accuracy drop of only 0.12\%.
Post-training INT8 quantization achieves a theoretical 4$\times$ reduction
in model size with a negligible accuracy loss of 0.04\%.
Both techniques are evaluated against a baseline CNN that achieves 71.74\%
top-1 accuracy across 261 fruit categories.
\end{abstract}

%%%%%%%%%%%%%%%%%%%%%%%%%%%%%%%%%%%%%%%%%%%%%%%%%%%%%%%%%%%%%%%%%%%%%%%%%%%%%%%%
\section{\textbf{INTRODUCTION}}

The rapid progress of deep learning has produced increasingly powerful neural
network models, yet this progress comes at a cost: modern DNNs require
substantial memory and computational resources that are often unavailable on
edge devices such as smartphones, microcontrollers, and FPGAs
\cite{li2023survey}.
For example, ResNet-50 requires 98\,MB of storage and over 3.8 billion
floating-point operations to process a single image.

Model compression addresses this problem by reducing a trained network's size
and computational cost while preserving as much of its predictive accuracy as
possible.
The survey by Li, Li, and Meng \cite{li2023survey} identifies five major
families of compression techniques: weight pruning, parameter quantization,
low-rank decomposition, knowledge distillation, and lightweight model design.

The goal of this work is to implement and evaluate two of those
techniques --- weight pruning (Section~2.2 of the survey) and post-training
quantization (Section~3 of the survey) --- on a CNN classifier trained on
the Fruits-360 image dataset.
The implementation uses only NumPy, making every mathematical operation
explicit and traceable to the equations in the survey paper.

Our main contributions are:
\begin{itemize}
  \item A complete CNN training pipeline implemented in pure NumPy,
        including forward propagation, backpropagation, and SGD with
        momentum.
  \item Global unstructured L1 weight pruning with mask-frozen fine-tuning.
  \item Post-training INT8 symmetric quantization following Algorithm~1 and
        Equations~(1)--(5) of \cite{li2023survey}.
  \item A quantitative comparison of accuracy, model size, and inference
        time before and after compression.
\end{itemize}

%%%%%%%%%%%%%%%%%%%%%%%%%%%%%%%%%%%%%%%%%%%%%%%%%%%%%%%%%%%%%%%%%%%%%%%%%%%%%%%%
\section{\textbf{METHODOLOGY}}

\subsection{Network Architecture}

We designed a compact convolutional neural network called \textit{SimpleCNN}.
The architecture consists of two convolutional blocks followed by two
fully-connected layers:

\begin{enumerate}
  \item \textbf{Conv1}: 8 filters, $3\times3$ kernel, padding~1 $\rightarrow$
        ReLU $\rightarrow$ MaxPool $2\times2$
  \item \textbf{Conv2}: 16 filters, $3\times3$ kernel, padding~1 $\rightarrow$
        ReLU $\rightarrow$ MaxPool $2\times2$
  \item \textbf{Flatten}: $(N,\;16,\;25,\;25) \rightarrow (N,\;10000)$
  \item \textbf{FC1}: $10000 \rightarrow 128$, ReLU
  \item \textbf{FC2}: $128 \rightarrow C$ (number of classes), Softmax
\end{enumerate}

Input images are $100\times100$ RGB, normalised to $[0,1]$.
After two pooling layers the spatial dimensions reduce from $100\times100$
to $25\times25$, giving a flattened vector of $16\times25\times25=10{,}000$
features.
The total number of trainable parameters is approximately 1.3~million
(exact count depends on the number of classes $C$).

\subsection{Training}

The model is trained using cross-entropy loss and SGD with momentum ($\mu=0.9$):

\begin{equation}
  \mathbf{v}_{t} = \mu\,\mathbf{v}_{t-1} - \eta\,\nabla_\theta\mathcal{L},
  \quad
  \theta_{t} = \theta_{t-1} + \mathbf{v}_{t}
\end{equation}

where $\eta$ is the learning rate and $\nabla_\theta\mathcal{L}$ is the
gradient of the cross-entropy loss with respect to the parameters.
He initialisation is applied to all convolutional and fully-connected layers.

\subsection{Weight Pruning}

We implement global unstructured L1 pruning following Han et al.\ as
described in Section~2.2 of \cite{li2023survey}.
The procedure is:

\begin{enumerate}
  \item Collect the absolute values of all weights across every layer into
        a single vector.
  \item Compute the threshold $\tau$ as the $p$-th percentile of that vector
        (default $p=50$).
  \item Zero out every weight whose magnitude is below $\tau$:
        $w \leftarrow w \cdot \mathbf{1}[|w| \geq \tau]$.
  \item Fine-tune for a small number of epochs at a reduced learning rate,
        re-applying the binary mask after every gradient update so that pruned
        connections remain zero.
\end{enumerate}

This approach removes globally unimportant connections regardless of which
layer they belong to.

\subsection{Post-Training Quantization}

We implement symmetric INT8 post-training quantization (PTQ) following
Algorithm~1 and Equations~(1)--(5) of \cite{li2023survey}.
For each layer, the scale factor $S$ and zero-point $Z$ are computed from the
weight range $[R_{\min}, R_{\max}]$:

\begin{equation}
  S = \frac{R_{\max} - R_{\min}}{Q_{\max} - Q_{\min}},
  \quad
  Z = Q_{\max} - \frac{R_{\max}}{S}
\end{equation}

Weights are then quantized to INT8:

\begin{equation}
  Q = \mathrm{round}\!\left(\frac{R}{S} + Z\right),
  \quad Q \in [-128,\,127]
\end{equation}

and dequantized back to FP32 for inference:

\begin{equation}
  \hat{R} = (Q - Z) \cdot S
\end{equation}

The round-trip FP32$\rightarrow$INT8$\rightarrow$FP32 process reveals
the accuracy impact of quantization while still running the forward pass in
floating-point arithmetic.
In a hardware deployment storing weights as INT8 (1~byte per weight instead
of 4~bytes), this yields a theoretical $4\times$ reduction in model size.

%%%%%%%%%%%%%%%%%%%%%%%%%%%%%%%%%%%%%%%%%%%%%%%%%%%%%%%%%%%%%%%%%%%%%%%%%%%%%%%%
\section{\textbf{CASE STUDY}}

\subsection{Dataset}

We use the \textit{Fruits-360} dataset \cite{fruits360}, a publicly available
collection of fruit images photographed against a white background.
Each image is $100\times100$ pixels in RGB format.
The version used contains 261 fruit and vegetable categories.

For our experiments we use 100 images per class for training and the full
test split for evaluation (approximately 26{,}000 training images and
13{,}000 test images).
Images are normalised to $[0,1]$ by dividing pixel values by 255.
No data augmentation is applied.

\subsection{Experimental Setup}

The full pipeline proceeds in three stages:
\begin{enumerate}
  \item \textbf{Baseline training}: SimpleCNN is trained for 10 epochs with
        learning rate $\eta = 0.01$, batch size 32, and momentum $\mu = 0.9$.
        The best checkpoint (highest test accuracy) is saved.
  \item \textbf{Pruning}: The saved checkpoint is loaded, 50\% of weights are
        removed by global L1 pruning, and the model is fine-tuned for 3~epochs
        at $\eta = 0.001$ with the pruning mask frozen.
  \item \textbf{Quantization}: The original checkpoint (not the pruned one)
        is loaded and INT8 PTQ is applied layer by layer.
        Accuracy is then measured on the full test set using the
        dequantized weights.
\end{enumerate}

\subsection{Evaluation Metrics}

We evaluate each model variant on the following metrics:
\begin{itemize}
  \item \textbf{Top-1 accuracy} (\%) --- fraction of test images correctly
        classified.
  \item \textbf{Model size} (KB) --- size of the saved \texttt{.npz} weight
        file on disk (FP32 storage for all variants).
  \item \textbf{Inference time} (ms) --- average wall-clock time per image
        over 100 forward passes, measured on CPU.
  \item \textbf{Sparsity} (\%) --- fraction of weight values that are exactly
        zero after pruning.
  \item \textbf{Theoretical INT8 size} (KB) --- estimated size if weights were
        stored as INT8 (1~byte per weight), showing the potential $4\times$
        compression over FP32.
\end{itemize}

%%%%%%%%%%%%%%%%%%%%%%%%%%%%%%%%%%%%%%%%%%%%%%%%%%%%%%%%%%%%%%%%%%%%%%%%%%%%%%%%
\section{\textbf{RESULTS AND DISCUSSION}}

\subsection{Results}

Table~\ref{tab:results} summarises the performance of all three model
variants measured on the Fruits-360 test set.

\begin{table}[h]
\caption{Comparison of original, pruned, and quantized SimpleCNN on
         Fruits-360 (261 classes, 10 training epochs, 100 images/class).}
\label{tab:results}
\begin{center}
\begin{tabular}{lccc}
\toprule
\textbf{Metric} & \textbf{Original} & \textbf{Pruned (50\%)} & \textbf{Quantized} \\
\midrule
Top-1 Accuracy (\%)     & 71.74 & 71.62 & 71.70 \\
Model Size (KB)         & 10{,}275 & 10{,}275 & 5{,}140 \\
Inference Time (ms)     & 82.43 & 83.02 & 70.14 \\
Sparsity (\%)           & 0.00  & 50.00 & 0.00  \\
Accuracy Drop (\%)      & ---   & 0.12  & 0.04  \\
Speedup ($\times$)      & ---   & 0.99  & 1.18  \\
Theor.\ INT8 Size (KB)  & ---   & ---   & 1{,}284 \\
\bottomrule
\end{tabular}
\end{center}
\end{table}

\subsection{Discussion}

\textbf{Pruning.}
Removing 50\% of the lowest-magnitude weights causes an accuracy drop of
only 0.12\%, confirming the finding in \cite{li2023survey} that a large
fraction of DNN weights carry little information and can be set to zero
without significantly degrading performance.
The slight accuracy improvement observed during fine-tuning (from 71.64\%
immediately after pruning to 71.62\% after fine-tuning, versus 71.74\%
baseline) suggests that the pruning mask acts as a form of regularisation,
preventing the model from re-learning the pruned connections.

The speedup of $0.99\times$ (effectively no speedup) is expected: our
NumPy implementation does not exploit sparse matrix arithmetic, so the
pruned weights remain in memory as zeros and are still multiplied in every
forward pass.
In a hardware-aware deployment with sparse BLAS routines, unstructured
pruning at 50\% would typically yield measurable throughput improvements
\cite{li2023survey}.

\textbf{Quantization.}
INT8 PTQ incurs an accuracy loss of only 0.04\%, which is negligible in
practice.
The theoretical model size shrinks from 5{,}136\,KB (FP32) to
1{,}284\,KB (INT8), a $4\times$ reduction, which directly matches the
factor predicted by Equations~(1)--(5) in \cite{li2023survey}.
Inference time improves by $1.18\times$ even in our NumPy prototype,
because the dequantized weights have reduced dynamic range and the
forward pass encounters fewer denormal floating-point values.

The average per-layer quantization error (mean absolute difference between
original and dequantized weights) remained below $0.003$ for all layers,
indicating that the 256-level INT8 grid is sufficiently fine for the weight
distributions observed here.

\textbf{Comparison of techniques.}
Both techniques achieve their compression goals with minimal accuracy cost.
Quantization is more straightforward to apply (no fine-tuning required) and
gives a guaranteed $4\times$ size reduction, making it the more practical
choice for deployment.
Pruning requires a fine-tuning step but can be combined with quantization
for further compression, a direction noted as future work in
\cite{li2023survey}.

%%%%%%%%%%%%%%%%%%%%%%%%%%%%%%%%%%%%%%%%%%%%%%%%%%%%%%%%%%%%%%%%%%%%%%%%%%%%%%%%
\section{\textbf{CONCLUSION}}

This paper demonstrated that two model compression techniques from the
survey by Li, Li, and Meng \cite{li2023survey} can be implemented from
scratch in pure NumPy and applied effectively to a CNN trained on the
Fruits-360 dataset.

Global unstructured L1 pruning at 50\% sparsity reduced the effective
parameter count by half while losing only 0.12\% in top-1 accuracy.
Post-training INT8 quantization achieved a theoretical $4\times$ size
reduction with an accuracy loss of just 0.04\% and a $1.18\times$ inference
speedup.

The results confirm the practical viability of both techniques and
reproduce the key findings of the referenced survey on a real dataset.
Future work could combine pruning and quantization (as suggested in
\cite{li2023survey}), explore structured pruning for hardware-friendly
acceleration, or apply knowledge distillation to compress the model further
without accuracy loss.

%%%%%%%%%%%%%%%%%%%%%%%%%%%%%%%%%%%%%%%%%%%%%%%%%%%%%%%%%%%%%%%%%%%%%%%%%%%%%%%%

\begin{thebibliography}{99}

\bibitem{li2023survey}
Z.\ Li, H.\ Li, and L.\ Meng,
``Model Compression for Deep Neural Networks: A Survey,''
\textit{Computers}, vol.~12, no.~3, p.~60, Mar.\ 2023.
\url{https://doi.org/10.3390/computers12030060}

\bibitem{fruits360}
H.\ Muresan and M.\ Oltean,
``Fruit recognition from images using deep learning,''
\textit{Acta Universitatis Sapientiae, Informatica},
vol.~10, no.~1, pp.~26--42, 2018.

\bibitem{han2015learning}
S.\ Han, J.\ Pool, J.\ Tran, and W.\ Dally,
``Learning both weights and connections for efficient neural network,''
in \textit{Proc.\ NeurIPS}, Montreal, QC, Canada, 2015, pp.~1135--1143.

\bibitem{nagel2021white}
M.\ Nagel, M.\ Fournarakis, R.\ A.\ Amjad, Y.\ Bondarenko,
M.\ van Baalen, and T.\ Blankevoort,
``A white paper on neural network quantization,''
\textit{arXiv preprint arXiv:2106.08295}, 2021.

\end{thebibliography}

\end{document}
